\documentclass[11pt]{article}
\usepackage[margin=1in]{geometry}
\usepackage{amsmath,amssymb}
\usepackage{booktabs}
\usepackage{bbm}
\newcommand{\ind}{\mathbbm{1}}
\newcommand{\RCI}{\mathrm{RCI}}
\title{Decomposing Changes in the Rotational Complexity Index}
\author{}
\date{}

\begin{document}
\maketitle

\section{The rotational complexity index}

For a field observed over a six-year window, let $c=(c_1,\dots,c_6)$ be the
sequence of crop codes. Following Socolar et al.\ (2021), the rotational
complexity index is built from four components:
\begin{align}
n(c)   &= \#\{\text{distinct crops in } c\} && \text{(richness)},\\
T_1(c) &= \textstyle\sum_{t=1}^{5}\ind[c_{t+1}\neq c_t] && \text{(one-year turnover)},\\
T_2(c) &= \textstyle\sum_{t=1}^{4}\ind[c_{t+2}\neq c_t] && \text{(two-year turnover)},\\
T_p(c) &= \textstyle\sum_{t=2}^{6}\ind[c_t=c_{t-1}]\,\ind[c_t\in\mathcal{P}]
          && \text{(perennial pseudo-turnover)},
\end{align}
where $\mathcal{P}$ is the set of perennial codes (alfalfa, pasture, etc.).
The index is
\begin{equation}
\RCI(c)=\sqrt{\,n(c)\cdot S(c)\,},
\qquad
S(c)=\tfrac{1}{2}\bigl(T_1(c)+T_2(c)+T_p(c)\bigr).
\label{eq:rci}
\end{equation}
$T_p$ is a correction, not an independent penalty: it restores the turnover
credit that a repeated perennial forfeits in $T_1$ and $T_2$, so that keeping a
soil-building perennial in place is not scored as a loss of complexity. Note the
factor $\tfrac12$ inside $S$; it is required to reproduce Socolar's published
values (e.g.\ CSCSCS $=\sqrt5\approx2.24$, six distinct crops $=\sqrt{27}\approx5.2$)
and is sometimes dropped in informal restatements of the formula.

\section{Decomposing a change in RCI}

We attribute a change in complexity between a ``before'' window $A$ and an
``after'' window $B$ to its components. Summarize each window by the vector
$\theta=(n,T_1,T_2,T_p)$ and write $\RCI = g(\theta)$ with
$g(\theta)=\sqrt{\tfrac12\,\theta_n(\theta_{T_1}+\theta_{T_2}+\theta_{T_p})}$.
The target is
\begin{equation}
\Delta\RCI = g(\theta^B)-g(\theta^A).
\end{equation}

Because $g$ is a square root of a product, it is nonlinear in its inputs: a
change in RCI has no exact additive split into per-component pieces, and a naive
first-order (Taylor) attribution leaves a remainder that grows with the size of
the change. We therefore use the \emph{Shapley} decomposition, which is exact.
Treating the four components as ``players'' whose ``payoff'' is $g$, define the
characteristic function $v(Q)=g(\theta^{Q})$, where $\theta^{Q}$ takes the
$B$-value for components in $Q\subseteq K=\{n,T_1,T_2,T_p\}$ and the $A$-value for
the rest. The contribution of component $k$ is its average marginal effect over
all orderings in which components are switched from $A$ to $B$:
\begin{equation}
\phi_k=\!\!\sum_{Q\subseteq K\setminus\{k\}}
\frac{|Q|!\,(|K|-|Q|-1)!}{|K|!}\,
\bigl[v(Q\cup\{k\})-v(Q)\bigr].
\label{eq:shapley}
\end{equation}
The decisive property is \emph{efficiency}: the contributions sum to the total
change with no residual,
\begin{equation}
\sum_{k\in K}\phi_k=g(\theta^B)-g(\theta^A)=\Delta\RCI .
\label{eq:efficiency}
\end{equation}
Shapley values also satisfy symmetry (equal contributions for
interchangeable components), the null-player property (a component with
$\theta_k^A=\theta_k^B$ receives $\phi_k=0$), and additivity across fields, so
per-field contributions can be summed or averaged to give sample-level
attributions whose shares add to $100\%$.

\section{Two component groupings}

Because $T_p$ exists only to offset $T_1$ (and $T_2$), the two move in opposite
directions whenever a perennial run lengthens: $T_1$ falls, $T_p$ rises. Read
separately, the negative $T_1$ contribution can be misread as ``less
complexity.'' We therefore report two groupings, both exact by
\eqref{eq:efficiency}:
\begin{itemize}
\item[(i)] \textbf{Four-way}: $K=\{n,\,T_1,\,T_2,\,T_p\}$. Isolates the perennial
correction as its own channel.
\item[(ii)] \textbf{Merged three-way}: $K=\{n,\,\tilde T_1,\,T_2\}$ with effective
year-to-year turnover $\tilde T_1=T_1+T_p$. Folds the correction into the
turnover it corrects, giving the net (interpretable) turnover contribution. This
is the preferred view for perennial (alfalfa) regimes.
\end{itemize}
For purely annual transitions $T_p\equiv0$, so the two groupings coincide.

\section{Worked example and expected results}

Consider a field moving from a corn--soybean rotation to a five-year alfalfa
run followed by corn:
\[
A=\text{C\,S\,C\,S\,C\,S}\quad\longrightarrow\quad B=\text{A\,A\,A\,A\,A\,C}
\qquad(\text{A}=\text{alfalfa, perennial}).
\]
The components and index are:
\begin{center}
\begin{tabular}{lccccc}
\toprule
 & $n$ & $T_1$ & $T_2$ & $T_p$ & $\RCI$ \\
\midrule
$A$ (C S C S C S) & 2 & 5 & 0 & 0 & $\sqrt5\approx2.236$ \\
$B$ (A A A A A C) & 2 & 1 & 1 & 4 & $\sqrt6\approx2.449$ \\
\bottomrule
\end{tabular}
\end{center}
so $\Delta\RCI = 2.449-2.236 = +0.213$. The Shapley contributions under both
groupings are:
\begin{center}
\begin{tabular}{lcc}
\toprule
Component & Four-way $\phi$ & Merged $\phi$ \\
\midrule
$n$ (richness)              & $0.000$  & $0.000$ \\
$T_1$                       & $-0.949$ & --- \\
$T_p$ (perennial)           & $+0.925$ & --- \\
$\tilde T_1=T_1+T_p$        & ---      & $0.000$ \\
$T_2$                       & $+0.238$ & $+0.213$ \\
\midrule
$\sum\phi\;(=\Delta\RCI)$   & $+0.213$ & $+0.213$ \\
\bottomrule
\end{tabular}
\end{center}

\medskip
\noindent\textbf{Reading the result.} Richness is unchanged ($n:2\to2$), so
$\phi_n=0$ by the null-player property. In the four-way split the perennial run
produces a large negative $T_1$ ($-0.949$) almost exactly cancelled by a large
positive $T_p$ ($+0.925$); taken alone the $T_1$ term is misleading. The merged
view resolves this: effective year-to-year turnover is unchanged
($\tilde T_1:5\to5$, so $\phi_{\tilde T_1}=0$), and \emph{the entire increase in
complexity is carried by two-year turnover} ($+0.213$). This is the substantive
statement one wants to report. Note that regrouping is not the same as summing
columns: $\phi_{T_2}$ itself shifts ($0.238\!\to\!0.213$) and
$\phi_{T_1}+\phi_{T_p}=-0.024\neq\phi_{\tilde T_1}=0$, because the Shapley
allocation is recomputed over the reduced set of players rather than collapsed
from the four-way one.

\medskip
\noindent\textbf{What to expect in general.}
\begin{itemize}
\item \emph{Exact additivity.} The contributions reproduce $\Delta\RCI$ to
numerical precision ($|\sum_k\phi_k-\Delta\RCI|\sim10^{-12}$); this is a built-in
check, not a fitted quantity.
\item \emph{Perennial regimes.} Expect $\phi_{T_1}$ and $\phi_{T_p}$ to be large
and opposite in the four-way split; the merged $\tilde T_1$ contribution is the
honest net effect. In annual regimes the two groupings are identical.
\item \emph{Coupled channels.} $n$ and turnover are not independent levers: adding
a new species raises richness \emph{and} tends to raise turnover, so their
contributions should be interpreted jointly, not as separable causes. Shapley
allocates the interaction correctly, but the narrative should reflect the
coupling.
\item \emph{Aggregation.} By additivity, the share of the sample's total (or
mean) $\Delta\RCI$ attributable to each component is well defined and sums to one,
supporting statements such as ``$x\%$ of the average complexity gain was
two-year turnover.''
\end{itemize}

\section{Marginal (Taylor) decomposition}

An alternative attribution differentiates \eqref{eq:rci} and weights each
component change by its local sensitivity. Writing $S=\tfrac12(T_1+T_2+T_p)$, the
partial derivatives are
\begin{equation}
\frac{\partial\RCI}{\partial n}=\frac{S}{2\,\RCI}
=\frac{T_1+T_2+T_p}{4\,\RCI},
\qquad
\frac{\partial\RCI}{\partial T_1}
=\frac{\partial\RCI}{\partial T_2}
=\frac{\partial\RCI}{\partial T_p}
=\frac{n}{4\,\RCI},
\end{equation}
giving the first-order attribution
\begin{equation}
\Delta\RCI\;\approx\;
\frac{T_1+T_2+T_p}{4\,\RCI}\,\Delta n
\;+\;\frac{n}{4\,\RCI}\,(\Delta T_1+\Delta T_2+\Delta T_p),
\label{eq:taylor}
\end{equation}
with the gradient evaluated at a chosen base point (e.g.\ window $A$). The
derivatives are themselves informative: richness matters more where turnover is
already high, and turnover matters more where richness is high, so the same
$\Delta$ in a component moves RCI by different amounts depending on the rotation
it starts from.

Unlike \eqref{eq:shapley}, this split is \emph{not exact}. It carries a
second-order remainder that grows with the size of the change and depends on the
base point at which the gradient is evaluated. For the worked example above,
evaluating \eqref{eq:taylor} at window $A$ gives component pieces
$(\phi_n,\phi_{T_1},\phi_{T_2},\phi_{T_p})=(0.000,-0.894,+0.224,+0.894)$
summing to $+0.224$, which overshoots the true $\Delta\RCI=+0.213$ by about
$5\%$; the residual does not vanish and is not distributed to any component.
Because rotation components change in discrete, often large steps, this gap is
the rule rather than the exception.

The marginal decomposition is therefore the right tool for a \emph{sensitivity}
question---``around a typical rotation, which lever moves RCI most?''---but not
for attributing an observed change, where the residual leaves the pieces failing
to reconcile to $\Delta\RCI$. We use \eqref{eq:shapley} for attribution and
reserve \eqref{eq:taylor} for interpreting local responsiveness.

\section{Implementation notes}

The change is computed between windows for the same field. Two design choices
matter. First, the \emph{unit of change}: overlapping year-to-year windows share
five of six years, so their differences largely reflect the window boundary
sliding rather than a shift in the cropping system; non-overlapping era windows
(e.g.\ 2008--2013 vs.\ 2014--2019) make $\Delta\RCI$ a genuine regime change and
avoid the induced autocorrelation. Second, \emph{completeness}: any window
containing a non-agricultural year yields an undefined $\RCI$ and is dropped, so
the decomposition is run on fields with a complete window at both endpoints of
the comparison.

\begin{thebibliography}{9}
\bibitem{socolar2021}
Socolar, Y., Goldstein, B.~R., de Valpine, P., \& Bowles, T.~M. (2021).
Biophysical and policy factors predict simplified crop rotations in the US
Midwest. \emph{Environmental Research Letters}, 16(5), 054045.
\end{thebibliography}

\end{document}